\chapter{Problem Discussions}

\section{Assignment 1}

\begin{convention}
By a ring, we mean a commutative ring with $1$. Also, we exclude the zero ring.
\end{convention}

\begin{enumerate}[label=\arabic*.]

\item Let $F$ be a field. Prove that $F[X]$ has infinitely many irreducible polynomials.

\item Let $R$ be a ring. Show that $R[X]$ is never a local ring.

\item Let $R$ be a ring in which every element is either nilpotent or a unit. Prove that $R$ is a local ring. How many prime ideals can $R$ have?

\item Let $R$ be a ring. Prove that $a \in J(R)$ if and only if $1+ab$ is a unit for all $b \in R$. (In particular, $1+a$ is a unit.)

\item Let $R$ be a ring and $I$ be a finitely generated ideal such that $I^2=I$. Show that there exists $e \in I$ such that $I = \langle e \rangle$ where $e^2=e$ (such an element is called idempotent).

\item Let $R$ be a ring and $I$ be a finitely generated ideal such that $I = \langle a_1,\cdots,a_n \rangle + I^2$. Show that there exists $e \in I^2$ such that $I = \langle a_1,\ldots,a_n,e\rangle$ where $e(1-e) \in \langle a_1,\cdots,a_n\rangle$. What happens if $R$ is a local ring?

\item Is $R[X] \hookrightarrow R[X,X^{-1}]$ an integral extension?

\item Let $A \hookrightarrow B$ be an integral extension where $A,B$ are integral domains. Let $C$ be an integral domain and $\varphi: B \to C$ be a morphism such that $\varphi|_A$ is injective. Prove that $\varphi$ is injective.

\item Determine the integral closure of $\bbz[\sqrt3]$ in $\bbq(\sqrt3)$.

\item Determine the integral closure of $\bbz[\sqrt5]$ in $\bbq(\sqrt5)$.

\end{enumerate}

\section{Assignment 2}

\begin{convention}
By a ring, we mean a commutative ring with $1$. Also, we exclude the zero ring.
\end{convention}

\begin{enumerate}[label=\arabic*.]

\item Let $p \in \bbz$ be a prime and write $\mathfrak p = \langle p \rangle$. Describe the elements of the local ring $\bbz_{\mathfrak p}$. What is the field of fractions of $\bbz_{\mathfrak p}$?

\item If $N,P$ are submodules of an $R$-module $M$, prove that:
\begin{enumerate}[label=(\alph*)]
    \item $S^{-1}(N+P) = S^{-1}N + S^{-1}P$;
    \item $S^{-1}(N\cap P) = S^{-1}N \cap S^{-1}P$;
    \item $S^{-1}(M/N) \cong S^{-1}M/S^{-1}N$ as $S^{-1}R$-modules.
\end{enumerate}

\item Let $R$ be a ring and $S \subset R$ be a multiplicative subset. We have the canonical ring morphism $f: R \to S^{-1}R$ given by $f(r) = \frac r1$. For an ideal $I$ of $R$, $I^e$ denotes the ideal in $S^{-1}R$ generated by $f(I)$ (this is called \emph{extension} of $I$). If $J \subseteq S^{-1}R$ is an ideal, then $f^{-1}(J)$ is an ideal of $R$ which will be denoted as $J^c$ (called the \emph{contraction} of $J$).
\begin{enumerate}[label=(\alph*)]
    \item Prove that any ideal of $S^{-1}R$ is an extended ideal.
    \item Prove that an ideal $I$ of $R$ is contraction of an ideal of $S^{-1}R$ if and only if no element of $S$ is a zero-divisor in $R/I$.
    \item $\Spec(S^{-1}R)$ is in bijective correspondence with $X = \{\mathfrak p \in \Spec(R) \mid \mathfrak p \cap S = \emptyset\}$.
    \item For ideals $I,J$ of $R$, show that: $S^{-1}(I+J) = S^{-1}I + S^{-1}J$, $S^{-1}(I \cap J) = S^{-1}I \cap S^{-1}J$, $S^{-1}(IJ) = S^{-1}I\, S^{-1}J$, $S^{-1}(\sqrt I) = \sqrt{S^{-1}I}$.
    \item If $N$ is the nilradical of $R$ then prove that the nilradical of $S^{-1}R$ is $S^{-1}N$.
\end{enumerate}

\item Let $\mathfrak p \in \Spec(R)$. Then $R_{\mathfrak p}$ is a local ring with unique maximal ideal $\mathfrak p R_{\mathfrak p}$ (we proved this). Therefore, $\dfrac{R_{\mathfrak p}}{\mathfrak p R_{\mathfrak p}}$ is a field, called the \emph{residue field} at $\mathfrak p$. On the other hand, $R/\mathfrak p$ is an integral domain. Prove that its field of fractions is the same as $\dfrac{R_{\mathfrak p}}{\mathfrak p R_{\mathfrak p}}$.

\item Let $M$ be a finitely generated $R$-module and $S$ be a multiplicative subset of $R$. Prove that $S^{-1}(\Ann(M)) = \Ann(S^{-1}M)$.

\item Let $M$ be a finitely generated $R$-module and $S$ be a multiplicative subset of $R$. Prove that $S^{-1}M = 0$ if and only if there is $s \in S$ such that $sM=0$.

\item Let $J$ be an ideal of $R$. Show that $S := 1+J := \{1+a \mid a \in J\}$ is a multiplicative set. Prove that $S^{-1}J$ is contained in the Jacobson radical of $S^{-1}R$.

\item Let $R$ be a ring and $S,T$ be multiplicative subsets. Let $U$ be the image of $T$ in $S^{-1}R$. Define $ST := \{st \mid s \in S,\ t \in T\}$. Check that $ST$ is also a multiplicative subset. Prove that the rings $(ST)^{-1}R$ and $U^{-1}(S^{-1}R)$ are isomorphic.

\item Let $R$ be a ring such that $R_{\mathfrak p}$ has no nonzero nilpotent for every $\mathfrak p \in \Spec(R)$. Prove that $R$ has no nonzero nilpotent. If $R$ is a ring such that $R_{\mathfrak p}$ is an integral domain for every $\mathfrak p \in \Spec(R)$, then does it imply that $R$ is an integral domain?

\item Let $R$ be a (nonzero) ring and $X$ be the set of all multiplicative subsets $S$ of $R$ such that $0 \notin S$. Show that $X$ has a maximal element. Prove further that $T \in X$ is a maximal element of $X$ if and only if $R \setminus T$ is a minimal prime ideal of $R$.

\end{enumerate}

\section{Assignment 3}

\begin{convention}
By a ring, we mean a commutative ring with $1$. Also, we exclude the zero ring.
\end{convention}

\begin{enumerate}[label=\arabic*.]

\item Let $R$ be an integral domain. Prove that the following are equivalent:
\begin{enumerate}[label=(\alph*)]
    \item $R$ is integrally closed.
    \item $R_{\mathfrak p}$ is integrally closed for all $\mathfrak p \in \Spec(R)$.
    \item $R_{\mathfrak m}$ is integrally closed for all $\mathfrak m \in \Max(R)$.
\end{enumerate}

\item Let $R \hookrightarrow S$ be an integral extension.
\begin{enumerate}[label=(\alph*)]
    \item Let $a \in R$ be such that $a$ is a unit in $S$. Then show that $a$ is a unit in $R$.
    \item Prove that $\Jac(R) = \Jac(S) \cap R$.
\end{enumerate}

\item Let $R$ be a ring and $G$ be a finite subgroup of the group of automorphisms of $R$ (automorphism means ring isomorphism from $R$ to $R$, and they form a group with respect to composition). Let
\[
R^G := \{ r \in R \mid \sigma(r) = r\ \forall\, \sigma \in G \}.
\]
\begin{enumerate}[label=(\alph*)]
    \item Check that $R^G$ is a subring of $R$ and $R$ is integral over $R^G$.
    \item Let $P \in \Spec(R^G)$ and $X := \{Q \in \Spec(R) \mid Q \cap R^G = P\}$. If $Q_1,Q_2 \in X$, prove that there is some $\tau \in G$ such that $Q_1 = \tau(Q_2)$. Conclude that $X$ is finite.
\end{enumerate}

\item Let $R \hookrightarrow S$ be an integral extension and $I$ be an ideal of $R$. Prove that
\[
\sqrt{IS} = \{ s \in S \mid s^n + a_1s^{n-1} + \cdots + a_n = 0 \text{ for some } n \in \bbn \text{ and } a_i \in I\ \forall\, i \}.
\]

\item Let $R \hookrightarrow S$ be an integral extension and let $\mathfrak n \in \Max(B)$. Let $\mathfrak m = \mathfrak n \cap A$ (note that $\mathfrak m \in \Max(A)$). Is $B_{\mathfrak n}$ integral over $A_{\mathfrak m}$?

\item Let $R \hookrightarrow S$ be an extension of rings such that $B \setminus A$ is closed under multiplication. Prove that $A$ is integrally closed in $B$.

\item Let $R$ be a ring and $f \in R[X]$ be a nonconstant monic. Show that there is a ring $S$ that contains $R$ as a subring such that $f$ is a product of monic linear factors in $S[X]$.

\item Let $R \hookrightarrow S$ be an extension of rings. Let $f,g$ be two monic polynomials in $S[X]$. Let the coefficients of $fg$ are all integral over $R$. Prove that the all coefficients of $f$ and all the coefficients of $g$ are integral over $R$.

\item Let $R \hookrightarrow S$ be an extension of rings. Prove that $f \in S[X]$ is integral over $R[X]$ if and only if all the coefficients of $f$ are integral over $R$. Conclude that if $C$ is the integral closure of $R$ in $S$, then $C[X]$ is the integral closure of $R[X]$ in $S[X]$.

\item Let $A$ be an integrally closed domain with field of fractions $K$. Let $L$ be a finite Galois extension of $K$ with Galois group $G$. Let $B$ be the integral closure of $A$ in $L$. Show that $\sigma(B) = B$ for all $\sigma \in G$ and that $A = B^G = \{ b \in B \mid \sigma(b) = b\ \forall\, \sigma \in G \}$.

\end{enumerate}

\section{Assignment 4}

\begin{convention}
By a ring, we mean a commutative ring with $1$. Also, we exclude the zero ring.
\end{convention}

\begin{enumerate}[label=\arabic*.]

\item Show that $2 \otimes \overline 1$ is $0$ in $\bbz \otimes_{\bbz} \dfrac{\bbz}{2\bbz}$ but is not $0$ in $2\bbz \otimes_{\bbz} \dfrac{\bbz}{2\bbz}$.

\item Using the above exercise, or otherwise, prove that tensor product does not preserve exactness.

\item Let $a,b \in \bbz$ be coprime. Prove that $\dfrac{\bbz}{a\bbz} \otimes_{\bbz} \dfrac{\bbz}{b\bbz} = 0$. More generally, if the gcd of $a,b$ is $d$, prove that
\[
\frac{\bbz}{a\bbz} \otimes_{\bbz} \frac{\bbz}{b\bbz} \cong \frac{\bbz}{d\bbz}.
\]

\item Let $R$ be a ring and $I,J$ be ideals of $R$. Prove that $(R/I) \otimes_R (R/J) \cong R/(I+J)$.

\item Let $M$ be an $R$-module and $I \subseteq R$ be an ideal. Prove that $(R/I) \otimes_R M \cong M/IM$.

\item Prove that $\bbq/\bbz \otimes_{\bbz} \bbq/\bbz = 0$. Prove also that $\bbq \otimes_{\bbz} \bbq/\bbz = 0$.

\item Let $M$ be an $R$-module and $I \subseteq R$ be an ideal. Prove that there is an $R$-linear map $g: I \otimes_R M \to IM$ and that $g$ is surjective. Give an example to show that $g$ is not injective in general.

\item Let $R \hookrightarrow S$ be rings. Prove that $S \otimes_R R[X]$ is isomorphic to $S[X]$ as $R$-modules.

\item Let $R$ be an integral domain with quotient field $K$ and let $V$ be a $K$-vector space. Prove that $K \otimes_R V \cong V$ as $R$-modules.

\item Let $M,N,P$ be $R$-modules. Prove that $M \otimes (N\otimes P) \cong M\otimes N \otimes P \cong (M\otimes N)\otimes P$.

\end{enumerate}